% ==============================================================================
% annexes/templates.tex — UE SIS589
% Annexe C : Modèles Professionnels
% ==============================================================================

\chapter{Modèles Professionnels}
\label{ann:templates}

\begin{boxinfo}
Cette annexe fournit des modèles directement exploitables en situation
opérationnelle. Chaque modèle est précédé d'instructions d'utilisation.
Les champs entre crochets \texttt{[CHAMP]} sont à compléter.
\end{boxinfo}

% ==============================================================================
\section{Modèle~1 — Procès-Verbal de Saisie Informatique}
\label{ann:tpl-pv}
% ==============================================================================

\begin{boxinfo}
\textbf{Usage~:} À remplir à la main ou sur ordinateur lors de toute
acquisition forensique. Conserver l'original dans le dossier physique~;
scanner et attacher au rapport numérique. Ce modèle est conforme aux
exigences ISO~27037 et RFC~3227.
\end{boxinfo}

\medskip
\begin{center}
{\large\bfseries\color{CouleurPrimaire}
PROCÈS-VERBAL DE COLLECTE ET D'ACQUISITION FORENSIQUE}\\[0.5ex]
{\small Référence~: \texttt{PV-[ANNEE]-[MOIS]-[NUMERO]}}
\end{center}

\bigskip

\textbf{1. IDENTIFICATION DE LA MISSION}

\begin{table}[H]\centering\small
\begin{tabular}{p{5cm}p{8.5cm}}
\toprule
Référence de l'incident & \texttt{INC-[ANNEE]-[MOIS][JOUR]-[NNN]} \\
\midrule
Organisation & \rule{8.5cm}{0.4pt} \\
Mandat légal & \rule{8.5cm}{0.4pt} \\
Date et heure (UTC) & \rule{8.5cm}{0.4pt} \\
Lieu d'intervention & \rule{8.5cm}{0.4pt} \\
\bottomrule
\end{tabular}
\end{table}

\textbf{2. IDENTIFICATION DE L'INVESTIGATEUR}

\begin{table}[H]\centering\small
\begin{tabular}{p{5cm}p{8.5cm}}
\toprule
Nom et prénom & \rule{8.5cm}{0.4pt} \\
\midrule
Qualification & \rule{8.5cm}{0.4pt} \\
Organisation & \rule{8.5cm}{0.4pt} \\
Contact & \rule{8.5cm}{0.4pt} \\
Témoin présent & \rule{8.5cm}{0.4pt} \\
\bottomrule
\end{tabular}
\end{table}

\textbf{3. DESCRIPTION DU SYSTÈME SOURCE}

\begin{table}[H]\centering\small
\begin{tabular}{p{5cm}p{8.5cm}}
\toprule
Type de matériel & ☐~Serveur \quad ☐~Poste \quad ☐~Laptop \quad ☐~Mobile \quad ☐~Autre \\
\midrule
Marque / Modèle & \rule{8.5cm}{0.4pt} \\
Numéro de série & \rule{8.5cm}{0.4pt} \\
Système d'exploitation & \rule{8.5cm}{0.4pt} \\
Adresse IP / MAC & \rule{8.5cm}{0.4pt} \\
État lors de la saisie & ☐~Allumé \quad ☐~Éteint \quad ☐~Hibernation \\
Présence de chiffrement & ☐~Oui (BitLocker / VeraCrypt / Autre) \quad ☐~Non \\
\bottomrule
\end{tabular}
\end{table}

\textbf{4. SUPPORT DE DESTINATION}

\begin{table}[H]\centering\small
\begin{tabular}{p{5cm}p{8.5cm}}
\toprule
Type & ☐~Disque externe \quad ☐~NAS \quad ☐~Serveur réseau \\
\midrule
Marque / Modèle & \rule{8.5cm}{0.4pt} \\
Numéro de série & \rule{8.5cm}{0.4pt} \\
Chiffrement du support & ☐~Oui (VeraCrypt / BitLocker) \quad ☐~Non \\
\bottomrule
\end{tabular}
\end{table}

\textbf{5. MÉTHODE ET OUTILLAGE}

\begin{table}[H]\centering\small
\begin{tabular}{p{5cm}p{8.5cm}}
\toprule
Type d'acquisition & ☐~Image disque bit à bit \quad ☐~Dump mémoire \quad ☐~Capture réseau \\
\midrule
Bloqueur en écriture & ☐~Matériel (marque/SN~: \rule{4cm}{0.4pt}) \quad ☐~Logiciel \quad ☐~Absent \\
Outil utilisé & ☐~dd \quad ☐~dcfldd \quad ☐~FTK Imager \quad ☐~avml \quad ☐~Autre \\
Version de l'outil & \rule{8.5cm}{0.4pt} \\
Format de l'image & ☐~RAW (.dd) \quad ☐~E01 \quad ☐~AFF4 \quad ☐~LIME \\
Fichier(s) produit(s) & \rule{8.5cm}{0.4pt} \\
\bottomrule
\end{tabular}
\end{table}

\textbf{6. INTÉGRITÉ CRYPTOGRAPHIQUE}

\begin{table}[H]\centering\small
\begin{tabular}{p{3cm}p{11cm}}
\toprule
Algorithme & ☐~SHA-256 \quad ☐~SHA-512 \quad ☐~MD5 (déconseillé seul) \\
\midrule
Hash source & \texttt{\rule{11cm}{0.4pt}} \\
Hash image & \texttt{\rule{11cm}{0.4pt}} \\
Vérification & ☐~\textbf{IDENTIQUES} (intégrité confirmée) \quad ☐~DIVERGENTS (alerter) \\
\bottomrule
\end{tabular}
\end{table}

\textbf{7. OBSERVATIONS ET INCIDENTS DE COLLECTE}

\vspace{3ex}
\rule{\linewidth}{0.4pt}\\[2ex]
\rule{\linewidth}{0.4pt}\\[2ex]
\rule{\linewidth}{0.4pt}

\textbf{8. SIGNATURES}

\begin{table}[H]\centering\small
\begin{tabular}{p{7cm}p{7cm}}
\toprule
\textbf{Investigateur} & \textbf{Témoin / Représentant de l'organisation} \\
\midrule
\vspace{2cm} & \vspace{2cm} \\
Signature~: \rule{4cm}{0.4pt} & Signature~: \rule{4cm}{0.4pt} \\
Date/heure (UTC)~: & Date/heure (UTC)~: \\
\bottomrule
\end{tabular}
\end{table}

\clearpage

% ==============================================================================
\section{Modèle~2 — Fiche de Gestion d'Incident}
\label{ann:tpl-incident}
% ==============================================================================

\begin{boxinfo}
\textbf{Usage~:} Ouvrir une fiche dès la première alerte qualifiée.
Mettre à jour en temps réel pendant la gestion. Joindre au rapport final.
\end{boxinfo}

\medskip
\begin{center}
{\large\bfseries\color{CouleurPrimaire}
FICHE DE GESTION D'INCIDENT DE SÉCURITÉ}
\end{center}

\textbf{IDENTIFICATION}

\begin{table}[H]\centering\small\rowcolors{2}{TableauPair}{TableauImpair}
\begin{tabular}{p{5cm}p{8.5cm}}
\toprule
N° de ticket & \texttt{INC-[ANNEE]-[MOIS][JOUR]-[NNN]} \\
Date/heure de détection (UTC) & \rule{8.5cm}{0.4pt} \\
Date/heure d'ouverture & \rule{8.5cm}{0.4pt} \\
Source de détection & ☐~SIEM \quad ☐~EDR \quad ☐~Utilisateur \quad ☐~Partenaire \quad ☐~CERT \\
Incident Manager & \rule{8.5cm}{0.4pt} \\
Analyste en charge & \rule{8.5cm}{0.4pt} \\
\bottomrule
\end{tabular}
\end{table}

\textbf{QUALIFICATION}

\begin{table}[H]\centering\small\rowcolors{2}{TableauPair}{TableauImpair}
\begin{tabular}{p{5cm}p{8.5cm}}
\toprule
Type d'incident & ☐~Brute force \quad ☐~Malware \quad ☐~Ransomware \quad ☐~Fuite \quad ☐~DDoS \quad ☐~Autre \\
Gravité & ☐~\textcolor{red}{Critique} \quad ☐~\textcolor{orange}{Haute} \quad ☐~\textcolor{yellow!60!black}{Moyenne} \quad ☐~\textcolor{green!50!black}{Basse} \\
Systèmes affectés & \rule{8.5cm}{0.4pt} \\
Données affectées & ☐~Oui (préciser~: \rule{4.5cm}{0.4pt}) \quad ☐~Non \quad ☐~Inconnu \\
IP source & \rule{8.5cm}{0.4pt} \\
Utilisateur(s) concerné(s) & \rule{8.5cm}{0.4pt} \\
TTP MITRE principaux & \rule{8.5cm}{0.4pt} \\
\bottomrule
\end{tabular}
\end{table}

\textbf{LOG DES ACTIONS (compléter en temps réel)}

\begin{table}[H]\centering\small
\begin{tabular}{p{3.5cm}p{2cm}p{7.5cm}p{0.8cm}}
\toprule
\tableauHeader{Timestamp (UTC)} & \tableauHeader{Acteur} &
\tableauHeader{Action effectuée} & \tableauHeader{OK} \\
\midrule
\rule{3.5cm}{0.4pt} & \rule{2cm}{0.4pt} & \rule{7.5cm}{0.4pt} & ☐ \\
\rule{3.5cm}{0.4pt} & \rule{2cm}{0.4pt} & \rule{7.5cm}{0.4pt} & ☐ \\
\rule{3.5cm}{0.4pt} & \rule{2cm}{0.4pt} & \rule{7.5cm}{0.4pt} & ☐ \\
\rule{3.5cm}{0.4pt} & \rule{2cm}{0.4pt} & \rule{7.5cm}{0.4pt} & ☐ \\
\rule{3.5cm}{0.4pt} & \rule{2cm}{0.4pt} & \rule{7.5cm}{0.4pt} & ☐ \\
\rule{3.5cm}{0.4pt} & \rule{2cm}{0.4pt} & \rule{7.5cm}{0.4pt} & ☐ \\
\rule{3.5cm}{0.4pt} & \rule{2cm}{0.4pt} & \rule{7.5cm}{0.4pt} & ☐ \\
\rule{3.5cm}{0.4pt} & \rule{2cm}{0.4pt} & \rule{7.5cm}{0.4pt} & ☐ \\
\bottomrule
\end{tabular}
\end{table}

\textbf{NOTIFICATIONS EFFECTUÉES}

\begin{table}[H]\centering\small\rowcolors{2}{TableauPair}{TableauImpair}
\begin{tabular}{p{3cm}p{3cm}p{2.5cm}p{5cm}}
\toprule
\tableauHeader{Destinataire} & \tableauHeader{Canal} &
\tableauHeader{Timestamp} & \tableauHeader{Contenu résumé} \\
\midrule
RSSI & ☐~Tél ☐~Email & \rule{2.5cm}{0.4pt} & \rule{5cm}{0.4pt} \\
DG / Direction & ☐~Tél ☐~Email & \rule{2.5cm}{0.4pt} & \rule{5cm}{0.4pt} \\
ANTIC & ☐~Email ☐~Courrier & \rule{2.5cm}{0.4pt} & \rule{5cm}{0.4pt} \\
Partenaires & ☐~Email & \rule{2.5cm}{0.4pt} & \rule{5cm}{0.4pt} \\
CERT / CSIRT ext. & ☐~Email & \rule{2.5cm}{0.4pt} & \rule{5cm}{0.4pt} \\
\bottomrule
\end{tabular}
\end{table}

\textbf{CLÔTURE DE L'INCIDENT}

\begin{table}[H]\centering\small\rowcolors{2}{TableauPair}{TableauImpair}
\begin{tabular}{p{5cm}p{8.5cm}}
\toprule
Date/heure de clôture (UTC) & \rule{8.5cm}{0.4pt} \\
Durée totale de l'incident & \rule{8.5cm}{0.4pt} \\
Cause racine identifiée & ☐~Oui \quad ☐~Non (préciser~: \rule{3.5cm}{0.4pt}) \\
RETEX programmé & ☐~Oui (date~: \rule{4cm}{0.4pt}) \quad ☐~Non \\
\bottomrule
\end{tabular}
\end{table}

\clearpage

% ==============================================================================
\section{Modèle~3 — Rapport d'Incident (Structure)}
\label{ann:tpl-rapport}
% ==============================================================================

\begin{boxinfo}
\textbf{Usage~:} Structure à respecter pour tout rapport d'incident
de niveau~\gravite{haute} ou \gravite{critique}. Deux versions à produire~:
technique (RSSI, SOC) et executive (Direction).
\end{boxinfo}

\medskip
\begin{center}
{\large\bfseries\color{CouleurPrimaire}
RAPPORT D'INCIDENT DE SÉCURITÉ}\\[0.3ex]
{\small\color{CouleurBordInfo} Version~: ☐~Confidentiel \quad
☐~Usage interne \quad ☐~TLP:AMBER}
\end{center}

\textbf{EN-TÊTE DU RAPPORT}

\begin{table}[H]\centering\small\rowcolors{2}{TableauPair}{TableauImpair}
\begin{tabular}{p{4cm}p{9.5cm}}
\toprule
Référence & \texttt{RPT-INC-[ANNEE]-[MOIS][JOUR]-[NNN]} \\
Organisation & \rule{9.5cm}{0.4pt} \\
Incident & \texttt{INC-[ANNEE]-[MOIS][JOUR]-[NNN]} \\
Auteur & \rule{9.5cm}{0.4pt} \\
Version & v1.0 (Brouillon) \quad v1.1 (Revue RSSI) \quad v2.0 (Final) \\
Classification & ☐~Public \quad ☐~Interne \quad ☐~Confidentiel \quad ☐~Secret \\
\bottomrule
\end{tabular}
\end{table}

\textbf{STRUCTURE TYPE}

\begin{enumerate}
  \item \textbf{Résumé exécutif} (1~page)
        \begin{itemize}
          \item Nature et date de l'incident
          \item Systèmes et données affectés
          \item Impact métier (disponibilité, confidentialité, finances)
          \item Mesures prises
          \item Statut actuel (en cours / résolu)
          \item Décision(s) requise(s) de la direction
        \end{itemize}

  \item \textbf{Chronologie de l'incident}
        \begin{itemize}
          \item Timeline annotée (tableau~: Timestamp UTC / Source / Événement / TTP MITRE)
          \item Vecteur d'accès initial
          \item Progression de l'attaque
          \item Moment et méthode de détection
        \end{itemize}

  \item \textbf{Analyse technique}
        \begin{itemize}
          \item Vecteur initial (CVE, credential, phishing…)
          \item Propagation et pivots
          \item Mapping complet MITRE ATT\&CK
          \item Outils utilisés par l'attaquant (si identifiés)
          \item Attribution (si possible~: groupe, pays, motivation)
        \end{itemize}

  \item \textbf{Preuves collectées}
        \begin{itemize}
          \item Tableau~: Pièce / Type / Hash SHA-256 / Date collecte / Outil
          \item État de la chaîne de custody
        \end{itemize}

  \item \textbf{Impact}
        \begin{itemize}
          \item Technique~: systèmes compromis, durée d'indisponibilité
          \item Données~: volumes exfiltrés ou chiffrés (estimés)
          \item Financier~: coût estimé (remédiation + pertes d'exploitation)
          \item Légal~: obligations de notification déclenchées
        \end{itemize}

  \item \textbf{Mesures de réponse}
        \begin{itemize}
          \item Confinement~: mesures, date, efficacité
          \item Éradication~: éléments supprimés, comptes réinitialisés
          \item Rétablissement~: systèmes restaurés, date de reprise
        \end{itemize}

  \item \textbf{Recommandations}
        Tableau~: Action / Priorité / Responsable / Délai / KPI de vérification

  \item \textbf{Annexes}
        Logs bruts, captures d'écran, règles SIEM déclenchées,
        liste des IoC, chaîne de custody complète
\end{enumerate}

\clearpage

% ==============================================================================
\section{Modèle~4 — Notification d'Incident à l'ANTIC}
\label{ann:tpl-antic}
% ==============================================================================

\begin{boxinfo}
\textbf{Usage~:} À envoyer à l'ANTIC dès qu'un incident de sécurité
significatif est confirmé, conformément à la Loi~2010/012 du Cameroun.
Envoyer par email recommandé avec accusé de réception à
\texttt{incident@antic.cm} (adresse indicative~; vérifier l'adresse
officielle en vigueur).
\end{boxinfo}

\medskip
\begin{center}
{\bfseries\color{CouleurPrimaire}
NOTIFICATION D'INCIDENT DE SÉCURITÉ INFORMATIQUE}\\[0.3ex]
{\small Conformément à la Loi n°~2010/012 du 21~décembre~2010}
\end{center}

\bigskip

\textbf{À l'attention de~:} Monsieur le Directeur Général de l'ANTIC\\
\textbf{Objet~:} Notification d'incident de sécurité: \texttt{[NOM ENTREPRISE]}

\bigskip

\textbf{1. IDENTIFICATION DE L'ORGANISATION NOTIFIANTE}
\vspace{0.5ex}

\begin{table}[H]\centering\small\rowcolors{2}{TableauPair}{TableauImpair}
\begin{tabular}{p{5cm}p{8.5cm}}
\toprule
Dénomination sociale & \rule{8.5cm}{0.4pt} \\
Secteur d'activité & \rule{8.5cm}{0.4pt} \\
Adresse & \rule{8.5cm}{0.4pt} \\
Responsable déclarant & \rule{8.5cm}{0.4pt} \\
Fonction & \rule{8.5cm}{0.4pt} \\
Contact & \rule{8.5cm}{0.4pt} \\
\bottomrule
\end{tabular}
\end{table}

\textbf{2. DESCRIPTION DE L'INCIDENT}

\begin{table}[H]\centering\small\rowcolors{2}{TableauPair}{TableauImpair}
\begin{tabular}{p{5cm}p{8.5cm}}
\toprule
Nature de l'incident & ☐~Intrusion \quad ☐~Ransomware \quad ☐~Fuite données \quad ☐~DDoS \quad ☐~Autre \\
Date/heure détectée (UTC) & \rule{8.5cm}{0.4pt} \\
Date/heure probable de début & \rule{8.5cm}{0.4pt} \\
Systèmes affectés & \rule{8.5cm}{0.4pt} \\
Données personnelles concernées & ☐~Oui (nombre estimé~: \rule{3cm}{0.4pt}) \quad ☐~Non \\
Impact sur les services & \rule{8.5cm}{0.4pt} \\
\bottomrule
\end{tabular}
\end{table}

\textbf{3. MESURES PRISES}

\vspace{1.5ex}
\rule{\linewidth}{0.4pt}\\[1.5ex]
\rule{\linewidth}{0.4pt}\\[1.5ex]
\rule{\linewidth}{0.4pt}

\textbf{4. ASSISTANCE SOUHAITÉE}

☐~Investigation technique \qquad
☐~Coordination avec d'autres organismes \qquad
☐~Information uniquement

\bigskip

Fait à \rule{4cm}{0.4pt}, le \rule{3cm}{0.4pt}

\vspace{2ex}
Signature du responsable déclarant~: \rule{6cm}{0.4pt}

\clearpage

% ==============================================================================
\section{Modèle~5 — Plan de Remédiation}
\label{ann:tpl-remediation}
% ==============================================================================

\begin{boxinfo}
\textbf{Usage~:} À compléter lors du RETEX post-incident.
Valider par le RSSI, présenter à la direction, suivre mensuellement.
\end{boxinfo}

\medskip
\begin{center}
{\large\bfseries\color{CouleurPrimaire}
PLAN DE REMÉDIATION POST-INCIDENT}\\[0.3ex]
{\small Incident~: \texttt{INC-[REF]} \quad Date~: \rule{4cm}{0.4pt}
\quad Responsable~: \rule{4cm}{0.4pt}}
\end{center}

\begin{table}[H]
\centering\small
\caption{Plan de remédiation: à compléter}
\label{tab:tpl-remediation}
\rowcolors{2}{TableauPair}{TableauImpair}
\begin{tabular}{p{0.5cm}p{4cm}p{1.8cm}p{1.5cm}p{2cm}p{2cm}p{2cm}}
\toprule
\tableauHeader{\#} & \tableauHeader{Action} &
\tableauHeader{Impact} & \tableauHeader{Effort} &
\tableauHeader{Responsable} & \tableauHeader{Délai} &
\tableauHeader{KPI} \\
\midrule
1 & \rule{4cm}{0.4pt} & ☐~C ☐~H ☐~M & ☐~F ☐~M ☐~E & \rule{2cm}{0.4pt} & J+\rule{0.8cm}{0.4pt} & \rule{2cm}{0.4pt} \\
2 & \rule{4cm}{0.4pt} & ☐~C ☐~H ☐~M & ☐~F ☐~M ☐~E & \rule{2cm}{0.4pt} & J+\rule{0.8cm}{0.4pt} & \rule{2cm}{0.4pt} \\
3 & \rule{4cm}{0.4pt} & ☐~C ☐~H ☐~M & ☐~F ☐~M ☐~E & \rule{2cm}{0.4pt} & J+\rule{0.8cm}{0.4pt} & \rule{2cm}{0.4pt} \\
4 & \rule{4cm}{0.4pt} & ☐~C ☐~H ☐~M & ☐~F ☐~M ☐~E & \rule{2cm}{0.4pt} & J+\rule{0.8cm}{0.4pt} & \rule{2cm}{0.4pt} \\
5 & \rule{4cm}{0.4pt} & ☐~C ☐~H ☐~M & ☐~F ☐~M ☐~E & \rule{2cm}{0.4pt} & J+\rule{0.8cm}{0.4pt} & \rule{2cm}{0.4pt} \\
6 & \rule{4cm}{0.4pt} & ☐~C ☐~H ☐~M & ☐~F ☐~M ☐~E & \rule{2cm}{0.4pt} & J+\rule{0.8cm}{0.4pt} & \rule{2cm}{0.4pt} \\
\bottomrule
\multicolumn{7}{l}{\small\textit{Impact~: C=Critique, H=Haute, M=Moyenne.
  Effort~: F=Faible, M=Moyen, E=Élevé.}} \\
\end{tabular}
\end{table}
